% ============================================================================
% CAPÍTULO 4 — EFS4AC
% Arquivo preparado para inclusão no projeto principal da dissertação.
%
% Observações:
% 1) O capítulo foi estruturado como um capítulo metodológico de dissertação,
%    apresentando a proposta, a arquitetura em etapas, as decisões de projeto
%    e a integração com a classificação associativa.
% 2) Mantém o label cap:Metodologia para não quebrar referências anteriores e
%    adiciona um label específico do método.
% 3) Pacotes recomendados no preâmbulo do projeto: amsmath, amssymb, array.
% ============================================================================
\chapter{\textit{EFS4AC}: Método de \textit{Ensemble Feature Selection} via Agregação de Rankings para Classificadores Associativos}\label{cap:Metodologia}

Este capítulo apresenta o método proposto nesta dissertação, denominado
\textit{EFS4AC}. A proposta foi concebida para incorporar \textit{Ensemble
Feature Selection} (EFS) ao contexto da Classificação Associativa, de modo a
selecionar um subconjunto de atributos antes da indução das regras de
classificação. A ideia central consiste em combinar diferentes métodos de
\textit{Feature Selection} (FS), transformar suas saídas em rankings
comparáveis, analisar a complementaridade entre eles, agregar essas evidências
em um consenso e, por fim, definir o subconjunto final de atributos a ser
utilizado pelo classificador associativo.

A motivação do \textit{EFS4AC} decorre de três constatações principais.
Primeiramente, a literatura mostra que diferentes métodos de FS capturam
critérios distintos de relevância, redundância e interação entre atributos,
de forma que a escolha de um único seletor pode introduzir dependência excessiva
de uma única perspectiva de avaliação \citeonline{bolon2019ensembles,salman2026review}.
Em segundo lugar, no contexto dos classificadores associativos, o número de
atributos influencia diretamente o espaço de busca das regras, afetando tanto o
custo computacional de indução quanto a complexidade do modelo final. Por fim,
a revisão realizada neste trabalho indicou uma lacuna na interseção entre
\textit{Ensemble Feature Selection} e Classificação Associativa, o que justifica
a proposição de um método específico para esse cenário.

Diferentemente de abordagens que alteram diretamente o mecanismo interno do
classificador associativo, o \textit{EFS4AC} atua como uma etapa anterior à
indução do modelo. Assim, o método não modifica a estrutura básica do
classificador associativo adotado; em vez disso, ele controla quais atributos
serão disponibilizados à mineração das regras. Essa decisão metodológica torna a
proposta modular, reprodutível e comparável tanto com a ausência de seleção de
atributos quanto com métodos individuais de FS.

% ============================================================================
\section{Visão Geral do Método}
\label{sec:ac-efs-visao-geral}
% ============================================================================

O fluxo geral do \textit{EFS4AC} é composto por cinco etapas principais:

\begin{enumerate}
    \item seleção dos métodos-base de \textit{Feature Selection};
    \item geração e padronização dos rankings individuais;
    \item análise e controle da diversidade entre os rankings;
    \item agregação dos rankings em um consenso global;
    \item definição do subconjunto final de atributos.
\end{enumerate}

A Figura~\ref{fig:fluxo-geral-ac-efs} apresenta essas cinco etapas na forma de
um fluxo metodológico. Cada caixa representa um bloco funcional do método, e as
setas indicam a dependência entre as saídas de uma etapa e as entradas da etapa
seguinte.

\begin{figure}[ht]
    \centering
    \small
    \setlength{\tabcolsep}{2pt}
    \renewcommand{\arraystretch}{1.3}
    \begin{tabular}{c@{\;$\rightarrow$\;}c@{\;$\rightarrow$\;}c@{\;$\rightarrow$\;}c@{\;$\rightarrow$\;}c}
        \fbox{\parbox[c][1.8cm][c]{2.5cm}{\centering \textbf{Etapa 1}\\ Seleção dos\\ métodos-base de FS}}
        &
        \fbox{\parbox[c][1.8cm][c]{2.45cm}{\centering \textbf{Etapa 2}\\ Geração e\\ padronização\\ dos rankings}}
        &
        \fbox{\parbox[c][1.8cm][c]{2.3cm}{\centering \textbf{Etapa 3}\\ Análise da\\ diversidade\\ dos rankings}}
        &
        \fbox{\parbox[c][1.8cm][c]{2.3cm}{\centering \textbf{Etapa 4}\\ Agregação\\ dos rankings}}
        &
        \fbox{\parbox[c][1.8cm][c]{2.4cm}{\centering \textbf{Etapa 5}\\ Definição do\\ subconjunto\\ final}}
    \end{tabular}
    \caption{Fluxo geral do método \textit{EFS4AC}.}
    \label{fig:fluxo-geral-ac-efs}
\end{figure}

Embora a Figura~\ref{fig:fluxo-geral-ac-efs} sintetize a proposta de forma
compacta, cada etapa envolve decisões metodológicas específicas. Por esse
motivo, a Tabela~\ref{tab:descricao-etapas-ac-efs} descreve resumidamente o
papel de cada caixa apresentada na figura.

\begin{table}[ht]
    \centering
    \small
    \caption{Descrição resumida das cinco etapas do método \textit{EFS4AC}.}
    \label{tab:descricao-etapas-ac-efs}
    \begin{tabular}{p{1.3cm}p{3.5cm}p{8.8cm}}
        \hline
        \textbf{Etapa} & \textbf{Bloco} & \textbf{Descrição} \\
        \hline
        1 & Seleção dos métodos-base de FS & Define quais seletores irão compor o ensemble, buscando compatibilidade com dados discretizados e diversidade de critérios de avaliação. \\
        2 & Geração e padronização dos rankings & Aplica cada seletor ao conjunto de treinamento e converte suas saídas para uma representação comum: um ranking completo dos atributos. \\
        3 & Análise da diversidade dos rankings & Mede a similaridade entre os rankings produzidos, permitindo identificar redundância, complementaridade e diferentes estratégias de composição do ensemble. \\
        4 & Agregação dos rankings & Combina os rankings individuais para produzir um ranking consensual global, utilizando um agregador de preferência. \\
        5 & Definição do subconjunto final & Converte o ranking consensual em um subconjunto efetivo de atributos, que será então utilizado pelo classificador associativo. \\
        \hline
    \end{tabular}
\end{table}

Formalmente, seja $D^{tr}=(X,Y)$ o conjunto de treinamento, em que
$X=\{f_1,f_2,\ldots,f_n\}$ representa o conjunto de atributos discretizados e
$Y$ representa a variável classe. Seja ainda
$\mathcal{S}=\{S_1,S_2,\ldots,S_M\}$ o conjunto de $M$ métodos-base de FS que
compõem o ensemble. Cada seletor $S_j$ produz um ranking individual $R_j$ sobre
os $n$ atributos. O conjunto desses rankings é então processado pelas etapas de
diversidade e agregação, gerando um ranking consensual $R^{*}$. Por fim, uma
função de corte $\mathcal{K}(\cdot)$ transforma esse ranking em um subconjunto
final $F^{*}\subseteq X$, conforme expresso em

\begin{equation}
    F^{*}=\mathcal{K}(R^{*}).
    \label{eq:subconjunto-final-ac-efs}
\end{equation}

Esse subconjunto $F^{*}$ é a saída principal do método e constitui a entrada da
etapa de indução do classificador associativo.

% ============================================================================
% ============================================================================
% ============================================================================
\section{Etapa 1 -- Seleção dos Métodos-Base de \textit{Feature Selection}}
\label{sec:ac-efs-etapa1}
% ============================================================================

A primeira etapa do método \textit{EFS4AC} consiste na definição dos
métodos-base de \textit{Feature Selection} (FS) que irão compor o ensemble.
Essa definição constitui uma decisão metodológica importante, uma vez que os
seletores escolhidos determinam quais perspectivas sobre a relevância dos
atributos estarão disponíveis para as etapas posteriores de geração, análise e
agregação dos rankings.

A utilização de múltiplos métodos de FS parte do princípio de que diferentes
seletores podem avaliar os mesmos atributos segundo critérios distintos e,
consequentemente, produzir resultados diferentes. No contexto de
\textit{Ensemble Feature Selection} (EFS), essa característica pode ser
explorada por meio de ensembles heterogêneos, nos quais diferentes métodos de
seleção são combinados com o objetivo de incorporar múltiplas perspectivas
sobre os dados. Dessa forma, busca-se reduzir a dependência de um único
critério de relevância e fornecer uma base mais diversificada para a construção
do consenso entre os seletores.

No \textit{EFS4AC}, optou-se por uma configuração heterogênea formada por
métodos supervisionados da categoria \textit{filter}. Além de não dependerem
diretamente do algoritmo de classificação utilizado posteriormente, esses
métodos permitem avaliar os atributos a partir de propriedades estatísticas e
informacionais dos dados. A utilização de filtros também favorece a construção
de uma representação comum baseada em rankings, necessária para as etapas
posteriores do método.

Do ponto de vista metodológico, esta etapa busca responder à seguinte questão:
\textit{quais métodos devem compor o ensemble de seleção de atributos?}
A resposta a essa questão é orientada por critérios relacionados tanto à
compatibilidade dos seletores com o contexto da Classificação Associativa
quanto à arquitetura proposta para o \textit{EFS4AC}. A
Figura~\ref{fig:fluxo-etapa1-ac-efs} apresenta o fluxo geral dessa etapa.

\begin{figure}[ht]
    \centering
    \small
    \setlength{\tabcolsep}{2pt}
    \renewcommand{\arraystretch}{1.3}

    \begin{tabular}{c@{\;$\rightarrow$\;}c@{\;$\rightarrow$\;}c}

        \fbox{\parbox[c][1.9cm][c]{3.0cm}{
            \centering
            \textbf{Entrada}\\
            Conjunto candidato de\\
            métodos de FS
        }}

        &

        \fbox{\parbox[c][1.9cm][c]{5.0cm}{
            \centering
            \textbf{Critérios de inclusão}\\
            dados discretizados; classe nominal;\\
            ranking completo; diversidade\\
            dos critérios de avaliação
        }}

        &

        \fbox{\parbox[c][1.9cm][c]{4.0cm}{
            \centering
            \textbf{Saída}\\
            Métodos-base elegíveis e\\
            composições do ensemble
        }}

    \end{tabular}

    \caption{Fluxo da Etapa 1 do método \textit{EFS4AC}.}
    \label{fig:fluxo-etapa1-ac-efs}
\end{figure}


% ----------------------------------------------------------------------------
\subsection{Critérios de inclusão dos métodos}
\label{subsec:ac-efs-criterios-inclusao}
% ----------------------------------------------------------------------------

A seleção dos métodos-base do \textit{EFS4AC} é orientada por quatro critérios
principais:

\begin{itemize}

    \item \textbf{compatibilidade com dados discretizados}: os métodos devem
    ser aplicáveis à representação discretizada utilizada pela Classificação
    Associativa, evitando a necessidade de reconstrução dos atributos contínuos
    originais;

    \item \textbf{aplicabilidade a problemas com classe nominal}: os métodos
    devem ser supervisionados e capazes de operar em problemas de classificação
    com variável-alvo categórica, contemplando cenários binários e multiclasse;

    \item \textbf{capacidade de produzir uma ordenação completa dos atributos}:
    cada método deve permitir, direta ou indiretamente, ordenar todos os
    atributos avaliados, uma vez que o \textit{EFS4AC} utiliza rankings como
    representação comum das saídas individuais;

    \item \textbf{diversidade dos critérios de avaliação}: o conjunto de
    métodos deve contemplar diferentes formas de avaliar a importância dos
    atributos, incluindo associação estatística, relevância informacional,
    redundância e interação ou complementaridade entre atributos.

\end{itemize}

Os três primeiros critérios estabelecem requisitos de compatibilidade com a
arquitetura do \textit{EFS4AC}, enquanto o último orienta a heterogeneidade do
ensemble. Dessa forma, a inclusão de um método não depende apenas de sua
aplicabilidade ao problema, mas também da perspectiva adicional que ele pode
incorporar ao conjunto de seletores.

É importante observar que a diversidade considerada nesta etapa é
\textit{conceitual}, isto é, decorre dos diferentes critérios utilizados pelos
métodos-base. A diversidade efetivamente observada entre os rankings produzidos
será analisada posteriormente, a partir da comparação entre as ordenações
geradas pelos seletores.


% ----------------------------------------------------------------------------
\subsection{Configuração inicial dos métodos-base}
\label{subsec:ac-efs-config-inicial}
% ----------------------------------------------------------------------------

Com base nos critérios estabelecidos, a configuração inicial do
\textit{EFS4AC} é composta por dez métodos: Gini, $\chi^2$,
\textit{Symmetric Uncertainty} (SU), \textit{Mutual Information Maximization}
(MIM), MRMR, JMI, CIFE, CMIM, ICAP e MRI.

A seleção desses métodos permite reunir duas perspectivas principais. A
primeira corresponde aos métodos de avaliação individual, nos quais a
relevância de cada atributo pode ser determinada diretamente em relação à
classe. A segunda corresponde aos métodos multivariados ou sequenciais, nos
quais a avaliação de um atributo considera, em diferentes graus, relações de
redundância, dependência condicional ou interação com atributos previamente
selecionados.

A Tabela~\ref{tab:metodos-base-ac-efs} apresenta os métodos-base e sintetiza
a contribuição conceitual de cada um para a composição do ensemble.

\begin{table}[ht]
    \centering
    \small
    \caption{Métodos-base da configuração inicial do \textit{EFS4AC}.}
    \label{tab:metodos-base-ac-efs}

    \begin{tabular}{p{2.0cm} p{3.0cm} p{7.2cm}}
        \hline
        \textbf{Método} &
        \textbf{Tipo} &
        \textbf{Critério predominante} \\
        \hline

        Gini
        & Univariado
        & Discriminação das classes a partir da distribuição dos dados. \\

        $\chi^2$
        & Univariado
        & Associação estatística entre o atributo e a classe. \\

        SU
        & Univariado
        & Dependência informacional normalizada entre atributo e classe. \\

        MIM
        & Univariado
        & Relevância individual baseada em informação mútua com a classe. \\

        \hline

        MRMR
        & Multivariado
        & Maximização da relevância e redução da redundância. \\

        JMI
        & Multivariado
        & Relevância conjunta e informação complementar entre atributos. \\

        CIFE
        & Multivariado
        & Relevância condicionada e interação com atributos selecionados. \\

        CMIM
        & Multivariado
        & Informação mútua condicional em relação aos atributos selecionados. \\

        ICAP
        & Multivariado
        & Relevância, redundância e efeitos de interação entre atributos. \\

        MRI
        & Multivariado
        & Relevância considerando redundância e interação informacional. \\

        \hline
    \end{tabular}
\end{table}

Os métodos Gini, $\chi^2$, SU e MIM constituem o grupo univariado da
configuração inicial. Embora avaliem os atributos individualmente, eles não
representam exatamente o mesmo critério: Gini e $\chi^2$ exploram propriedades
estatísticas ou distributivas, enquanto SU e MIM utilizam medidas derivadas da
teoria da informação.

MRMR, JMI, CIFE, CMIM, ICAP e MRI, por sua vez, constituem o grupo
multivariado ou sequencial. Esses métodos introduzem diferentes mecanismos para
considerar relações entre atributos durante a seleção, permitindo incorporar
aspectos como redundância, dependência condicional e complementaridade. A
literatura de FS baseada em teoria da informação contempla, por exemplo,
métodos consolidados como MRMR, JMI e CMIM, que diferem justamente na forma
como relevância e dependências entre atributos são consideradas.

Assim, a heterogeneidade da configuração inicial não decorre apenas da
quantidade de métodos empregados, mas principalmente da combinação de
diferentes critérios de avaliação. Essa característica é particularmente
relevante para o \textit{EFS4AC}, uma vez que a proposta busca investigar se
o consenso entre diferentes perspectivas de relevância pode produzir uma
seleção de atributos mais adequada à Classificação Associativa.


% ----------------------------------------------------------------------------
\subsection{Métodos reservados para investigação posterior}
\label{subsec:ac-efs-metodos-posteriores}
% ----------------------------------------------------------------------------

Além dos dez métodos da configuração inicial, foram considerados MIFS, DCSF,
CFR, IWFS e DISR. Esses métodos não são descartados da pesquisa, mas são
reservados para extensões ou análises posteriores, pois sua inclusão na
configuração inicial introduziria decisões metodológicas adicionais.

Em particular, MIFS requer a definição de um hiperparâmetro adicional
($\beta$); DCSF, CFR e IWFS apresentam particularidades relacionadas à
inicialização de seus procedimentos sequenciais; e DISR demanda uma análise
adicional quanto à obtenção de uma ordenação completa dos atributos. A
incorporação imediata desses métodos exigiria, portanto, parametrizações ou
adaptações que acrescentariam novos fatores ao experimento principal.

A decisão de mantê-los fora da configuração inicial busca preservar uma
interface homogênea entre os métodos-base quanto à geração posterior dos
rankings, reduzindo a necessidade de adaptações específicas para determinados
seletores. Ao mesmo tempo, esses métodos permanecem como alternativas para
avaliar, em experimentos complementares, a expansão do conjunto de componentes
do ensemble.


% ----------------------------------------------------------------------------
\subsection{Composições experimentais do ensemble}
\label{subsec:ac-efs-composicoes-ensemble}
% ----------------------------------------------------------------------------

A partir dos métodos selecionados, são definidas três composições experimentais
para investigar a influência da natureza dos seletores sobre o comportamento
do \textit{EFS4AC}:

\begin{itemize}

    \item $E_U$: ensemble composto pelos métodos univariados
    $\{\text{Gini}, \chi^2, \text{SU}, \text{MIM}\}$;

    \item $E_M$: ensemble composto pelos métodos multivariados ou sequenciais
    $\{\text{MRMR}, \text{JMI}, \text{CIFE}, \text{CMIM},
    \text{ICAP}, \text{MRI}\}$;

    \item $E_{U+M}$: ensemble heterogêneo completo, formado pela união
    $E_U \cup E_M$.

\end{itemize}

A configuração $E_{U+M}$ constitui a configuração principal do método, pois
combina avaliações individuais com critérios capazes de considerar relações
entre atributos. As configurações $E_U$ e $E_M$ são utilizadas como condições
experimentais para investigar se a combinação dessas duas perspectivas produz
comportamento distinto daquele observado quando o ensemble é composto apenas
por métodos de uma mesma natureza.

Essa comparação também permite analisar se eventuais diferenças observadas
decorrem da natureza dos seletores empregados ou da complementaridade entre
perspectivas univariadas e multivariadas. Como $E_U$, $E_M$ e $E_{U+M}$
possuem números distintos de componentes, o efeito da quantidade de métodos
deverá ser considerado na interpretação dos resultados experimentais.

Em síntese, a Etapa 1 estabelece quais perspectivas de seleção de atributos
serão incorporadas ao \textit{EFS4AC}. Sua saída é constituída pelo conjunto
de métodos-base elegíveis e pelas respectivas composições experimentais. A
etapa seguinte utiliza esses métodos para produzir uma representação comum na
forma de rankings completos dos atributos.
% ============================================================================
% ============================================================================
\section{Etapa 2 -- Geração e Padronização dos Rankings}
\label{sec:ac-efs-etapa2}
% ============================================================================

A segunda etapa do \textit{EFS4AC} tem como objetivo transformar as saídas dos
métodos-base definidos na Etapa 1 em uma representação comum. Essa padronização
é necessária porque os seletores empregados possuem diferentes mecanismos de
avaliação: enquanto alguns atribuem um escore individual a cada atributo,
outros constroem a seleção de forma sequencial, considerando os atributos já
selecionados. No \textit{EFS4AC}, essas diferentes saídas são convertidas em
\textbf{rankings completos dos atributos}.

A utilização das posições dos rankings, em vez dos valores brutos produzidos
pelos seletores, evita a comparação direta entre medidas definidas em escalas
distintas. Dessa forma, independentemente do critério empregado pelo método de
FS, todos os seletores passam a representar a preferência sobre os atributos
por meio de posições ordinais, nas quais a posição 1 corresponde ao atributo
mais preferido.

A Figura~\ref{fig:geracao-rankings-ac-efs} sintetiza esse procedimento,
distinguindo as duas formas de obtenção dos rankings utilizadas no método.

\begin{figure}[ht]
    \centering

    \begin{tikzpicture}[
        every node/.style={
            font=\small,
            align=center
        },
        metodo/.style={
            draw,
            rectangle,
            text width=3.0cm,
            minimum height=1.65cm,
            inner sep=5pt
        },
        processo/.style={
            draw,
            rectangle,
            text width=2.8cm,
            minimum height=1.65cm,
            inner sep=5pt
        },
        ranking/.style={
            draw,
            rectangle,
            text width=2.9cm,
            minimum height=2.3cm,
            inner sep=5pt
        },
        matriz/.style={
            draw,
            rectangle,
            text width=3.0cm,
            minimum height=2.3cm,
            inner sep=5pt
        },
        seta/.style={
            thick,
            ->,
            >=stealth
        }
    ]

    % ============================================================
    % COLUNA 1 -- Métodos
    % ============================================================

    \node[metodo] (uni) {
        \textbf{Métodos univariados}\\[2mm]
        Gini, $\chi^2$, SU, MIM
    };

    \node[metodo, below=0.65cm of uni] (multi) {
        \textbf{Métodos multivariados}\\[2mm]
        MRMR, JMI, CIFE,\\
        CMIM, ICAP, MRI
    };


    % ============================================================
    % COLUNA 2 -- Transformação
    % ============================================================

    \node[processo, right=0.70cm of uni] (score) {
        \textbf{Ordenação}\\[2mm]
        escores individuais
    };

    \node[processo, right=0.70cm of multi] (seq) {
        \textbf{Seleção sequencial}\\[2mm]
        ordem de entrada
    };


    % ============================================================
    % COLUNA 3 -- Representação comum
    % Posicionada explicitamente à direita das duas transformações
    % ============================================================

    \node[ranking, right=1.25cm of score, yshift=-1.15cm] (rank) {
        \textbf{Ranking completo}\\[3mm]
        $R_j$\\[2mm]
        $r_j(f_i)\in\{1,\ldots,n\}$\\[2mm]
        $1=$ melhor posição
    };


    % ============================================================
    % COLUNA 4 -- Saída
    % ============================================================

    \node[matriz, right=0.75cm of rank] (matriz) {
        \textbf{Matriz de rankings}\\[3mm]
        $\mathbf{R}\in\mathbb{R}^{n\times M}$\\[2mm]
        atributos $\times$ seletores
    };


    % ============================================================
    % SETAS
    % ============================================================

    \draw[seta]
        (uni.east) -- (score.west);

    \draw[seta]
        (multi.east) -- (seq.west);

    % Univariados -> ranking
    \draw[seta]
        (score.east) -- ++(0.45,0)
        |- ([yshift=0.45cm]rank.west);

    % Multivariados -> ranking
    \draw[seta]
        (seq.east) -- ++(0.45,0)
        |- ([yshift=-0.45cm]rank.west);

    % Ranking -> matriz
    \draw[seta]
        (rank.east) -- (matriz.west);

    \end{tikzpicture}

    \caption{Geração e padronização dos rankings na Etapa 2 do
    \textit{EFS4AC}.}
    \label{fig:geracao-rankings-ac-efs}
\end{figure}

Seja $S_j$ um dos $M$ seletores e seja
$F=\{f_1,f_2,\ldots,f_n\}$ o conjunto de atributos. O ranking produzido pelo
seletor $S_j$ é representado por

\begin{equation}
    R_j =
    \left(
        r_j(f_1),
        r_j(f_2),
        \ldots,
        r_j(f_n)
    \right),
    \label{eq:ranking-individual-ac-efs}
\end{equation}

em que $r_j(f_i)$ representa a posição atribuída ao atributo $f_i$ pelo
seletor $S_j$. Em todos os rankings, posições menores representam maior
preferência, de modo que $r_j(f_i)=1$ identifica o atributo mais bem
posicionado pelo método.


% ----------------------------------------------------------------------------
\subsection{Rankings provenientes de métodos univariados}
\label{subsec:ac-efs-univariados}
% ----------------------------------------------------------------------------

Gini, $\chi^2$, SU e MIM atribuem um escore individual a cada atributo.
Considerando $s_j(f_i)$ o escore produzido pelo método $S_j$ para o atributo
$f_i$, o ranking é obtido pela ordenação dos atributos segundo o sentido de
preferência da respectiva medida.

De forma geral, para critérios em que valores maiores indicam maior relevância,

\begin{equation}
    r_j(f_i)
    =
    1+
    \sum_{\substack{f_k\in F\\k\neq i}}
    \mathbb{I}
    \left[
        s_j(f_k)>s_j(f_i)
    \right],
    \label{eq:ranking-score-ac-efs}
\end{equation}

em que $\mathbb{I}[\cdot]$ representa a função indicadora. Caso determinado
critério utilize o sentido inverso de preferência, sua orientação será
normalizada antes da geração do ranking, mantendo-se a convenção de que a
primeira posição corresponde sempre ao atributo mais preferido.

Eventuais empates entre os escores serão tratados por uma regra determinística
de desempate, aplicada igualmente a todos os conjuntos de dados e execuções.
Essa padronização evita que decisões aleatórias de implementação introduzam
variação artificial nos rankings produzidos.


% ----------------------------------------------------------------------------
\subsection{Rankings provenientes de métodos multivariados}
\label{subsec:ac-efs-multivariados}
% ----------------------------------------------------------------------------

MRMR, JMI, CIFE, CMIM, ICAP e MRI produzem a ordenação por meio de um
procedimento sequencial. Diferentemente dos métodos univariados, a preferência
por um atributo pode depender dos atributos selecionados nas iterações
anteriores.

Inicialmente, considera-se vazio o conjunto de atributos selecionados:

\begin{equation}
    A_0=\emptyset.
\end{equation}

Na iteração $t$, seleciona-se, entre os atributos ainda disponíveis, aquele
que maximiza o critério correspondente ao método $S_j$:

\begin{equation}
    f_t^{*}
    =
    \underset{f\in F\setminus A_{t-1}}
    {\arg\max}
    \;
    J_j(f\mid A_{t-1}),
    \label{eq:selecao-sequencial-ac-efs}
\end{equation}

em que $J_j(f\mid A_{t-1})$ representa a função de avaliação do seletor,
considerando o conjunto de atributos previamente selecionados.

Após a escolha de $f_t^{*}$, atualiza-se o conjunto:

\begin{equation}
    A_t=A_{t-1}\cup\{f_t^{*}\}.
    \label{eq:atualizacao-conjunto-ac-efs}
\end{equation}

A posição atribuída ao atributo selecionado na iteração $t$ é definida por

\begin{equation}
    r_j(f_t^{*})=t.
    \label{eq:rank-sequencial-ac-efs}
\end{equation}

O procedimento é repetido até $t=n$, de modo que todos os atributos sejam
incluídos. Assim, a ordem de inserção no conjunto $A_t$ determina diretamente
o ranking completo produzido pelo seletor.


% ----------------------------------------------------------------------------
\subsection{Padronização e saída da Etapa 2}
\label{subsec:ac-efs-saida-etapa2}
% ----------------------------------------------------------------------------

Independentemente do mecanismo utilizado para sua geração, todos os rankings
resultantes seguem a mesma convenção:

\begin{equation}
    1 \leq r_j(f_i) \leq n,
\end{equation}

sendo as menores posições associadas aos atributos mais preferidos. Dessa
forma, as diferenças de escala e de interpretação dos escores originais são
eliminadas antes das etapas de comparação e agregação.

Após a execução dos $M$ seletores, obtém-se o conjunto

\begin{equation}
    \mathcal{R}
    =
    \{R_1,R_2,\ldots,R_M\}.
\end{equation}

Os rankings são então organizados em uma matriz de postos
$\mathbf{R}\in\mathbb{R}^{n\times M}$:

\begin{equation}
    \mathbf{R}=
    \begin{bmatrix}
        r_1(f_1) & r_2(f_1) & \cdots & r_M(f_1) \\
        r_1(f_2) & r_2(f_2) & \cdots & r_M(f_2) \\
        \vdots   & \vdots   & \ddots & \vdots   \\
        r_1(f_n) & r_2(f_n) & \cdots & r_M(f_n)
    \end{bmatrix}.
    \label{eq:matriz-ranks-ac-efs}
\end{equation}

Nessa matriz, cada linha corresponde a um atributo e cada coluna ao ranking
produzido por um método-base. A matriz $\mathbf{R}$ constitui a saída da
Etapa 2 e fornece a representação comum utilizada na etapa seguinte para
analisar a similaridade e a diversidade entre os rankings.

% ============================================================================
% ============================================================================
\section{Etapa 3 -- Análise e Controle da Diversidade dos Rankings}
\label{sec:ac-efs-etapa3}
% ============================================================================

A terceira etapa do \textit{EFS4AC} tem como objetivo analisar o grau de
concordância entre os rankings produzidos pelos métodos-base e investigar se
rankers excessivamente semelhantes devem exercer influência independente na
agregação. Essa análise é realizada sobre a matriz de rankings $\mathbf{R}$
obtida na Etapa 2.

A motivação para essa etapa decorre da própria natureza de um ensemble
heterogêneo. Embora os métodos-base tenham sido selecionados por representarem
diferentes critérios de avaliação, essa diversidade conceitual não garante que
eles produzam rankings empiricamente distintos para um determinado conjunto de
dados. Dois seletores fundamentados em critérios diferentes podem resultar em
ordenações muito semelhantes e, nesse caso, sua participação simultânea na
agregação pode aumentar implicitamente a influência de uma mesma preferência.

Por outro lado, baixa concordância entre dois rankings não implica,
necessariamente, que um deles deva ser descartado. Um ranker divergente pode
capturar relações não identificadas pelos demais seletores. Dessa forma, o
\textit{EFS4AC} não assume que maior diversidade seja sempre desejável, mas
trata a similaridade entre rankings como uma propriedade a ser mensurada e
controlada experimentalmente.

A Figura~\ref{fig:diversidade-rankings-ac-efs} apresenta o fluxo geral da
Etapa 3.

\begin{figure}[ht]
    \centering

    \begin{tikzpicture}[
        every node/.style={
            font=\small,
            align=center
        },
        bloco/.style={
            draw,
            rectangle,
            text width=3.0cm,
            minimum height=1.65cm,
            inner sep=5pt
        },
        seta/.style={
            thick,
            ->,
            >=stealth
        }
    ]

    \node[bloco] (entrada) {
        \textbf{Entrada}\\[1mm]
        matriz de rankings\\
        $\mathbf{R}\in\mathbb{R}^{n\times M}$
    };

    \node[bloco, right=0.65cm of entrada] (similaridade) {
        \textbf{Similaridade}\\[1mm]
        correlação de\\
        Spearman $\rho_{ij}$
    };

    \node[bloco, right=0.65cm of similaridade] (matriz) {
        \textbf{Diagnóstico}\\[1mm]
        matriz de similaridade\\
        $\mathbf{S}$
    };

    \node[bloco, right=0.65cm of matriz] (controle) {
        \textbf{Controle}\\[1mm]
        manutenção, limiar\\
        ou agrupamento
    };

    \draw[seta] (entrada) -- (similaridade);
    \draw[seta] (similaridade) -- (matriz);
    \draw[seta] (matriz) -- (controle);

    \end{tikzpicture}

    \caption{Análise e controle da diversidade entre os rankings na
    Etapa 3 do \textit{EFS4AC}.}
    \label{fig:diversidade-rankings-ac-efs}
\end{figure}


% ----------------------------------------------------------------------------
\subsection{Medição da similaridade entre rankings}
\label{subsec:ac-efs-similaridade}
% ----------------------------------------------------------------------------

Sejam $R_i$ e $R_j$ dois rankings completos produzidos pelos seletores
$S_i$ e $S_j$. A concordância entre suas ordenações é estimada por meio do
coeficiente de correlação de postos de Spearman:

\begin{equation}
    \rho_{ij}
    =
    \rho(R_i,R_j),
    \qquad
    -1\leq\rho_{ij}\leq1.
    \label{eq:spearman-ac-efs}
\end{equation}

Valores de $\rho_{ij}$ próximos de 1 indicam que os dois seletores atribuem
posições semelhantes aos atributos. Valores próximos de zero indicam baixa
associação monotônica entre as ordenações, enquanto valores próximos de $-1$
representam preferências aproximadamente inversas.

A utilização da correlação de postos é compatível com a representação adotada
pelo \textit{EFS4AC}, pois a comparação é realizada diretamente sobre as
posições dos atributos, sem depender das escalas originais dos escores
produzidos pelos métodos-base.

As similaridades calculadas para todos os pares de seletores são organizadas
em uma matriz simétrica

\begin{equation}
    \mathbf{S}
    =
    \begin{bmatrix}
        1          & \rho_{12} & \cdots & \rho_{1M} \\
        \rho_{21}  & 1         & \cdots & \rho_{2M} \\
        \vdots     & \vdots    & \ddots & \vdots    \\
        \rho_{M1}  & \rho_{M2} & \cdots & 1
    \end{bmatrix},
    \label{eq:matriz-similaridade-ac-efs}
\end{equation}

em que cada elemento $S_{ij}=\rho_{ij}$ representa a concordância entre os
rankings produzidos por $S_i$ e $S_j$.

A matriz $\mathbf{S}$ fornece um diagnóstico da estrutura de concordância do
ensemble. Em particular, blocos de valores elevados indicam grupos de
seletores que tendem a produzir ordenações semelhantes e que, portanto, podem
introduzir redundância na etapa de agregação.


% ----------------------------------------------------------------------------
\subsection{Representação da dissimilaridade entre rankers}
\label{subsec:ac-efs-distancia}
% ----------------------------------------------------------------------------

Para estratégias que requerem uma representação de dissimilaridade, a
correlação é transformada em

\begin{equation}
    d_{ij}
    =
    1-\rho_{ij}.
    \label{eq:distancia-ac-efs}
\end{equation}

Nessa definição, rankings idênticos apresentam $d_{ij}=0$, enquanto a
dissimilaridade aumenta à medida que diminui sua concordância.

Optou-se por preservar o sinal da correlação, em vez de utilizar
$|\rho_{ij}|$. Essa decisão é importante porque rankings fortemente
anticorrelacionados não são considerados redundantes no \textit{EFS4AC}:
embora apresentem forte associação, eles representam preferências opostas sobre
os atributos. A redundância de interesse nesta etapa corresponde,
principalmente, à repetição de ordenações semelhantes.


% ----------------------------------------------------------------------------
\subsection{Estratégias de controle da redundância}
\label{subsec:ac-efs-estrategias-diversidade}
% ----------------------------------------------------------------------------

A influência do controle de diversidade será analisada por meio de três
configurações. A primeira mantém todos os rankings e funciona como referência
para verificar se a intervenção sobre a composição do ensemble é efetivamente
necessária. As demais utilizam a matriz de similaridade para reduzir a
influência de rankers redundantes.

\begin{enumerate}

    \item \textbf{Sem controle de redundância ($D_0$).}
    Todos os rankings produzidos na Etapa 2 são preservados:

    \begin{equation}
        \mathcal{R}_{D_0}=\mathcal{R}.
    \end{equation}

    Essa configuração constitui o cenário de referência e permite avaliar se
    os mecanismos de controle introduzidos nesta etapa proporcionam benefícios
    nas etapas posteriores.

    \item \textbf{Controle por limiar de similaridade ($D_L$).}
    Pares de rankings cuja concordância ultrapasse um limiar $\tau$ são
    identificados como potencialmente redundantes:

    \begin{equation}
        \rho_{ij}\geq\tau.
        \label{eq:limiar-similaridade-ac-efs}
    \end{equation}

    Nesses casos, evita-se que ambos contribuam de maneira independente para
    representar essencialmente a mesma ordenação. O valor de $\tau$ não será
    definido a partir dos resultados finais de classificação, devendo ser
    estabelecido previamente ou determinado exclusivamente nos dados de
    treinamento.

    \item \textbf{Controle por agrupamento de rankers ($D_C$).}
    A matriz de dissimilaridade $\mathbf{D}=[d_{ij}]$ é utilizada para
    identificar grupos de rankings semelhantes. Seja

    \begin{equation}
        \mathcal{C}
        =
        \{C_1,C_2,\ldots,C_K\}
    \end{equation}

    uma partição dos $M$ rankers em $K$ grupos. Cada grupo reúne métodos cujas
    ordenações apresentam maior concordância entre si. A contribuição dos
    grupos para a agregação é então controlada de modo a evitar que um conjunto
    numeroso de rankers muito semelhantes exerça influência desproporcional
    sobre o consenso.

\end{enumerate}

A comparação entre $D_0$, $D_L$ e $D_C$ permite avaliar empiricamente se o
controle explícito da redundância entre rankings contribui para o
\textit{EFS4AC}. Essa análise também evita assumir, \textit{a priori}, que a
remoção ou o agrupamento de rankings semelhantes necessariamente melhora o
ensemble.


% ----------------------------------------------------------------------------
\subsection{Saída da Etapa 3}
\label{subsec:ac-efs-saida-etapa3}
% ----------------------------------------------------------------------------

A saída da Etapa 3 depende da estratégia de controle considerada. No cenário
sem controle, a matriz original de rankings é preservada. Nas configurações
com controle de redundância, obtém-se uma estrutura refinada que identifica
quais rankings ou grupos de rankings participarão da etapa de agregação.

De forma geral, essa saída pode ser representada por

\begin{equation}
    \mathcal{R}^{*}
    =
    \{R^{*}_1,R^{*}_2,\ldots,R^{*}_{M'}\},
    \qquad
    M'\leq M,
    \label{eq:rankings-refinados-ac-efs}
\end{equation}

quando houver seleção de representantes, ou pela estrutura de grupos
$\mathcal{C}$ quando a informação de agrupamento for preservada para a
agregação.

Além dos rankings refinados, a matriz de similaridade $\mathbf{S}$ constitui
uma saída analítica da etapa, pois permite caracterizar o grau de concordância
existente entre os diferentes métodos-base.

Assim, a Etapa 3 atua como uma camada intermediária entre a geração dos
rankings e sua agregação. Seu objetivo não é maximizar a diversidade de forma
indiscriminada, mas identificar e controlar situações em que múltiplos
seletores representam essencialmente a mesma preferência, preservando, sempre
que possível, perspectivas complementares para a construção do ranking
consensual.
% ============================================================================
% ============================================================================
\section{Etapa 4 -- Agregação dos Rankings}
\label{sec:ac-efs-etapa4}
% ============================================================================

A quarta etapa do \textit{EFS4AC} tem como objetivo integrar os rankings
provenientes das etapas anteriores em uma única ordenação consensual dos
atributos. Enquanto a Etapa 3 analisa e, quando necessário, controla a
redundância entre os rankers, a presente etapa determina como as preferências
individuais serão combinadas para produzir uma decisão coletiva.

Seja

\begin{equation}
    \mathcal{R}^{*}
    =
    \{R^{*}_1,R^{*}_2,\ldots,R^{*}_{M'}\}
\end{equation}

o conjunto de rankings disponibilizado pela Etapa 3, em que $M'$ representa o
número de rankings efetivamente considerados na agregação. No caso em que
nenhum ranker tenha sido removido, tem-se $M'=M$.

Para cada atributo $f_i$, as posições

\begin{equation}
    \mathbf{r}(f_i)
    =
    \left(
        r_1(f_i),
        r_2(f_i),
        \ldots,
        r_{M'}(f_i)
    \right)
\end{equation}

representam as preferências atribuídas pelos diferentes seletores. A função de
agregação utiliza essas posições para calcular um escore consensual
$c(f_i)$:

\begin{equation}
    c(f_i)
    =
    \mathcal{A}
    \left(
        r_1(f_i),
        r_2(f_i),
        \ldots,
        r_{M'}(f_i)
    \right),
    \label{eq:agregacao-geral-ac-efs}
\end{equation}

em que $\mathcal{A}$ representa o operador de agregação adotado.

Os atributos são posteriormente ordenados de acordo com $c(f_i)$, produzindo
o ranking consensual

\begin{equation}
    R^{*}
    =
    \left(
        f_{(1)},f_{(2)},\ldots,f_{(n)}
    \right),
    \label{eq:ranking-consensual-ac-efs}
\end{equation}

em que $f_{(1)}$ representa o atributo de maior preferência segundo o
consenso do ensemble.

A Figura~\ref{fig:agregacao-ac-efs} sintetiza o procedimento da Etapa 4.

\begin{figure}[ht]
    \centering

    \begin{tikzpicture}[
        every node/.style={
            font=\small,
            align=center
        },
        bloco/.style={
            draw,
            rectangle,
            text width=3.0cm,
            minimum height=1.65cm,
            inner sep=5pt
        },
        agregador/.style={
            draw,
            rectangle,
            text width=3.3cm,
            minimum height=2.0cm,
            inner sep=5pt
        },
        seta/.style={
            thick,
            ->,
            >=stealth
        }
    ]

    \node[bloco] (entrada) {
        \textbf{Rankings de entrada}\\[1mm]
        $\mathcal{R}^{*}$\\
        provenientes da Etapa 3
    };

    \node[agregador, right=0.75cm of entrada] (agregacao) {
        \textbf{Agregação}\\[1mm]
        Borda\\
        RRF\\
        TOPSIS
    };

    \node[bloco, right=0.75cm of agregacao] (escores) {
        \textbf{Consenso}\\[1mm]
        escores agregados\\
        $c(f_i)$
    };

    \node[bloco, right=0.75cm of escores] (saida) {
        \textbf{Saída}\\[1mm]
        ranking consensual\\
        $R^{*}$
    };

    \draw[seta] (entrada) -- (agregacao);
    \draw[seta] (agregacao) -- (escores);
    \draw[seta] (escores) -- (saida);

    \end{tikzpicture}

    \caption{Agregação dos rankings e geração do ranking consensual na
    Etapa 4 do \textit{EFS4AC}.}
    \label{fig:agregacao-ac-efs}
\end{figure}


% ----------------------------------------------------------------------------
\subsection{Agregação por Borda}
\label{subsec:ac-efs-borda}
% ----------------------------------------------------------------------------

O método de Borda transforma a posição atribuída a cada atributo em pontos.
Considerando $n$ atributos, o atributo situado na primeira posição recebe
$n$ pontos, o segundo recebe $n-1$ e assim sucessivamente. Dessa forma, para
o atributo $f_i$, define-se

\begin{equation}
    c_{\mathrm{Borda}}(f_i)
    =
    \sum_{j=1}^{M'}
    \left(
        n-r_j(f_i)+1
    \right).
    \label{eq:borda-ac-efs}
\end{equation}

Quanto maior $c_{\mathrm{Borda}}(f_i)$, maior a preferência coletiva pelo
atributo. O ranking consensual é obtido pela ordenação decrescente desses
escores.

A utilização de Borda é particularmente adequada como estratégia de referência
por operar diretamente sobre posições e permitir que todos os rankings
contribuam para a decisão final. Entretanto, por realizar uma transformação
linear das posições, diferenças entre posições consecutivas exercem a mesma
variação na pontuação independentemente de ocorrerem no início ou no final do
ranking.


% ----------------------------------------------------------------------------
\subsection{Agregação por \textit{Reciprocal Rank Fusion}}
\label{subsec:ac-efs-rrf}
% ----------------------------------------------------------------------------

Como segunda estratégia, será investigado o \textit{Reciprocal Rank Fusion}
(RRF). Diferentemente da transformação linear utilizada por Borda, o RRF
atribui maior influência às primeiras posições dos rankings por meio de uma
função recíproca.

Para cada atributo $f_i$, define-se

\begin{equation}
    c_{\mathrm{RRF}}(f_i)
    =
    \sum_{j=1}^{M'}
    \frac{1}{k+r_j(f_i)},
    \label{eq:rrf-ac-efs}
\end{equation}

em que $k>0$ é uma constante de suavização. Valores maiores de
$c_{\mathrm{RRF}}(f_i)$ indicam maior consenso em favor do atributo.

O parâmetro $k$ controla a diferença relativa entre posições sucessivas.
Valores menores tornam o método mais sensível às primeiras posições, enquanto
valores maiores suavizam essas diferenças. Por essa razão, o valor utilizado
no experimento deverá ser previamente especificado e mantido constante entre
os conjuntos de dados, ou definido exclusivamente a partir dos dados de
treinamento, evitando sua escolha com base no conjunto de teste.


% ----------------------------------------------------------------------------
\subsection{Agregação por TOPSIS}
\label{subsec:ac-efs-topsis}
% ----------------------------------------------------------------------------

A terceira estratégia considera a agregação como um problema de decisão
multicritério por meio do \textit{Technique for Order Preference by Similarity
to Ideal Solution} (TOPSIS). Nessa formulação, cada atributo corresponde a uma
alternativa e cada ranker constitui um critério de decisão.

Como as posições originais seguem a convenção de que valores menores
representam maior preferência, elas são inicialmente convertidas em valores de
benefício:

\begin{equation}
    x_{ij}
    =
    n-r_j(f_i)+1,
    \label{eq:topsis-transformacao-ac-efs}
\end{equation}

de forma que valores maiores correspondam a melhores posições.

A partir desses valores é construída a matriz de decisão

\begin{equation}
    \mathbf{X}
    =
    \begin{bmatrix}
        x_{11} & x_{12} & \cdots & x_{1M'} \\
        x_{21} & x_{22} & \cdots & x_{2M'} \\
        \vdots & \vdots & \ddots & \vdots \\
        x_{n1} & x_{n2} & \cdots & x_{nM'}
    \end{bmatrix}.
    \label{eq:topsis-matriz-ac-efs}
\end{equation}

Em seguida, os valores são normalizados:

\begin{equation}
    z_{ij}
    =
    \frac{x_{ij}}
    {\sqrt{\sum_{i=1}^{n}x_{ij}^{2}}}.
    \label{eq:topsis-normalizacao-ac-efs}
\end{equation}

Considerando inicialmente contribuição uniforme dos rankers, atribui-se

\begin{equation}
    w_j=\frac{1}{M'},
    \qquad
    \sum_{j=1}^{M'}w_j=1,
    \label{eq:topsis-pesos-ac-efs}
\end{equation}

e obtém-se a matriz ponderada por

\begin{equation}
    v_{ij}=w_jz_{ij}.
    \label{eq:topsis-ponderacao-ac-efs}
\end{equation}

A solução ideal positiva e a solução ideal negativa são definidas,
respectivamente, por

\begin{equation}
    v_j^{+}
    =
    \max_i v_{ij},
    \qquad
    v_j^{-}
    =
    \min_i v_{ij}.
    \label{eq:topsis-ideais-ac-efs}
\end{equation}

Para cada atributo, calculam-se então suas distâncias em relação às duas
soluções:

\begin{equation}
    D_i^{+}
    =
    \sqrt{
        \sum_{j=1}^{M'}
        \left(v_{ij}-v_j^{+}\right)^2
    },
\end{equation}

\begin{equation}
    D_i^{-}
    =
    \sqrt{
        \sum_{j=1}^{M'}
        \left(v_{ij}-v_j^{-}\right)^2
    }.
\end{equation}

Finalmente, o coeficiente de proximidade do atributo à solução ideal é dado
por

\begin{equation}
    c_{\mathrm{TOPSIS}}(f_i)
    =
    \frac{D_i^{-}}
    {D_i^{+}+D_i^{-}}.
    \label{eq:topsis-score-ac-efs}
\end{equation}

Quanto maior $c_{\mathrm{TOPSIS}}(f_i)$, mais próximo o atributo se encontra
da solução ideal e, consequentemente, melhor sua posição no ranking consensual.


% ----------------------------------------------------------------------------
\subsection{Integração com a Etapa 3}
\label{subsec:ac-efs-agregacao-ponderada}
% ----------------------------------------------------------------------------

A formulação anterior considera, por padrão, que os rankers possuem a mesma
influência. Entretanto, quando a estratégia adotada na Etapa 3 produzir pesos
associados aos rankings ou aos grupos de rankings, esses pesos poderão ser
incorporados à agregação.

Seja

\begin{equation}
    \mathbf{w}
    =
    (w_1,w_2,\ldots,w_{M'}),
    \qquad
    w_j\geq0,
    \qquad
    \sum_{j=1}^{M'}w_j=1.
    \label{eq:pesos-rankers-ac-efs}
\end{equation}

Nesse caso, a versão ponderada de Borda pode ser expressa por

\begin{equation}
    c_{\mathrm{Borda}}^{(w)}(f_i)
    =
    \sum_{j=1}^{M'}
    w_j
    \left(n-r_j(f_i)+1\right),
    \label{eq:borda-ponderado-ac-efs}
\end{equation}

enquanto o RRF ponderado é definido por

\begin{equation}
    c_{\mathrm{RRF}}^{(w)}(f_i)
    =
    \sum_{j=1}^{M'}
    \frac{w_j}{k+r_j(f_i)}.
    \label{eq:rrf-ponderado-ac-efs}
\end{equation}

No TOPSIS, os mesmos pesos podem ser incorporados diretamente na matriz
normalizada ponderada por meio de $v_{ij}=w_jz_{ij}$.

Essa formulação permite que o mecanismo de controle de diversidade da Etapa 3
seja integrado à agregação sem alterar a estrutura dos agregadores. Quando
nenhuma ponderação for empregada, utiliza-se $w_j=1/M'$, recuperando-se a
configuração de influência uniforme.


% ----------------------------------------------------------------------------
\subsection{Comparação entre as estratégias de agregação}
\label{subsec:ac-efs-comparacao-agregadores}
% ----------------------------------------------------------------------------

Borda, RRF e TOPSIS representam mecanismos distintos de construção do
consenso. Borda utiliza uma transformação linear das posições; RRF enfatiza
não linearmente atributos posicionados nas primeiras colocações; e TOPSIS
interpreta os rankers como critérios de decisão e avalia cada atributo em
relação a soluções ideais.

A comparação entre essas estratégias será realizada sob as mesmas entradas e
o mesmo protocolo experimental. Dessa maneira, diferenças observadas nas
etapas posteriores poderão ser associadas ao mecanismo de agregação, sem
alterações simultâneas nos rankings fornecidos ao ensemble.


% ----------------------------------------------------------------------------
\subsection{Saída da Etapa 4}
\label{subsec:ac-efs-saida-etapa4}
% ----------------------------------------------------------------------------

Independentemente do agregador utilizado, a Etapa 4 produz dois elementos:
um escore consensual $c(f_i)$ para cada atributo e o correspondente ranking
global $R^{*}$.

Formalmente, após calcular os escores para todos os atributos,

\begin{equation}
    \mathbf{c}
    =
    \left(
        c(f_1),c(f_2),\ldots,c(f_n)
    \right),
\end{equation}

os atributos são ordenados do maior para o menor valor de consenso, resultando
em $R^{*}$.

A Etapa 4, portanto, não determina ainda quais atributos serão efetivamente
mantidos. Sua função é produzir uma preferência global a partir das evidências
fornecidas pelos diferentes métodos-base. A transformação desse ranking em um
subconjunto de atributos é realizada exclusivamente na Etapa 5, preservando a
separação entre os problemas de \textit{rank aggregation} e de definição do
ponto de corte.

% ============================================================================
% ============================================================================
\section{Etapa 5 -- Definição do Subconjunto Final de Atributos}
\label{sec:ac-efs-etapa5}
% ============================================================================

A quinta etapa do \textit{EFS4AC} transforma o ranking consensual produzido na
Etapa 4 em um subconjunto efetivo de atributos. Embora a agregação determine
uma ordenação global de preferência, ela não estabelece, por si só, quantos
atributos devem ser utilizados na indução do classificador associativo. Assim,
esta etapa responde à seguinte questão metodológica: \textit{em qual posição do
ranking consensual deve ser realizado o corte?}

Seja

\begin{equation}
    R^{*}
    =
    \left(
        f_{(1)},f_{(2)},\ldots,f_{(n)}
    \right)
\end{equation}

o ranking consensual, em que $f_{(1)}$ representa o atributo mais bem
posicionado, e seja

\begin{equation}
    \mathbf{c}^{*}
    =
    \left(
        c_{(1)},c_{(2)},\ldots,c_{(n)}
    \right)
\end{equation}

o vetor dos respectivos escores de consenso, ordenados de acordo com
$R^{*}$.

O subconjunto final é obtido por uma função de corte $\mathcal{K}$:

\begin{equation}
    F^{*}
    =
    \mathcal{K}
    \left(
        R^{*},\mathbf{c}^{*}
    \right),
    \qquad
    F^{*}\subseteq F.
    \label{eq:funcao-corte-ac-efs}
\end{equation}

A separação entre agregação e corte é deliberada. A Etapa 4 determina
\textit{quais atributos são preferidos pelo ensemble}, enquanto a Etapa 5
determina \textit{quantos desses atributos devem ser preservados}. Dessa forma,
diferentes estratégias de corte podem ser comparadas sobre um mesmo ranking
consensual sem modificar o mecanismo responsável por sua construção.

A Figura~\ref{fig:corte-ac-efs} apresenta o fluxo da Etapa 5.

\begin{figure}[ht]
    \centering

    \begin{tikzpicture}[
        every node/.style={
            font=\small,
            align=center
        },
        bloco/.style={
            draw,
            rectangle,
            text width=3.0cm,
            minimum height=1.7cm,
            inner sep=5pt
        },
        corte/.style={
            draw,
            rectangle,
            text width=3.5cm,
            minimum height=2.1cm,
            inner sep=5pt
        },
        seta/.style={
            thick,
            ->,
            >=stealth
        }
    ]

    \node[bloco] (entrada) {
        \textbf{Entrada}\\[1mm]
        ranking consensual $R^{*}$\\
        e escores $c(f_i)$
    };

    \node[corte, right=0.70cm of entrada] (regra) {
        \textbf{Regra de corte}\\[1mm]
        percentual fixo\\
        ou adaptativo
    };

    \node[bloco, right=0.70cm of regra] (k) {
        \textbf{Ponto de corte}\\[1mm]
        $k^{*}$ atributos
    };

    \node[bloco, right=0.70cm of k] (saida) {
        \textbf{Saída}\\[1mm]
        subconjunto final\\
        $F^{*}$
    };

    \draw[seta] (entrada) -- (regra);
    \draw[seta] (regra) -- (k);
    \draw[seta] (k) -- (saida);

    \end{tikzpicture}

    \caption{Definição do ponto de corte e obtenção do subconjunto final de
    atributos na Etapa 5 do \textit{EFS4AC}.}
    \label{fig:corte-ac-efs}
\end{figure}


% ----------------------------------------------------------------------------
\subsection{Cortes percentuais fixos}
\label{subsec:ac-efs-corte-percentual}
% ----------------------------------------------------------------------------

A primeira estratégia utiliza percentuais predefinidos do ranking consensual.
Essa configuração fornece uma referência controlada para avaliar o efeito da
quantidade de atributos sobre o classificador associativo.

Seja $p\in(0,1]$ a proporção de atributos a ser preservada. O número de
atributos selecionados é definido por

\begin{equation}
    k_p
    =
    \left\lceil
        p\,n
    \right\rceil.
    \label{eq:k-percentual-ac-efs}
\end{equation}

O subconjunto correspondente é

\begin{equation}
    F^{*}_{p}
    =
    \left\{
        f_{(1)},f_{(2)},\ldots,f_{(k_p)}
    \right\}.
    \label{eq:subset-percentual-ac-efs}
\end{equation}

Inicialmente, serão considerados os níveis

\begin{equation}
    p\in\{0{,}25,\;0{,}50,\;0{,}75\},
\end{equation}

correspondentes, respectivamente, à preservação dos primeiros 25\%, 50\% e
75\% dos atributos.

Esses cortes não pressupõem que determinado percentual seja universalmente
adequado. Sua função é fornecer condições experimentais padronizadas que
permitam observar o comportamento do \textit{EFS4AC} sob diferentes níveis
de redução dimensional.


% ----------------------------------------------------------------------------
\subsection{Corte adaptativo baseado no consenso}
\label{subsec:ac-efs-corte-adaptativo}
% ----------------------------------------------------------------------------

Além dos percentuais fixos, será investigada uma estratégia adaptativa capaz de
determinar o ponto de corte a partir da própria distribuição dos escores
produzidos pelo ensemble. A motivação é evitar que a quantidade de atributos
selecionados dependa exclusivamente de percentuais definidos previamente,
permitindo que o corte responda à estrutura de consenso observada em cada
conjunto de dados.

Como Borda, RRF e TOPSIS produzem escores em escalas distintas, os valores de
consenso são inicialmente transformados para uma escala comum. Seja
$c_{(i)}$ o escore do atributo situado na posição $i$ do ranking consensual.
Define-se o escore normalizado

\begin{equation}
    \tilde{c}_{(i)}
    =
    \frac{
        c_{(i)}-\min_{\ell}c_{(\ell)}
    }{
        \max_{\ell}c_{(\ell)}
        -
        \min_{\ell}c_{(\ell)}
    },
    \label{eq:normalizacao-consenso-ac-efs}
\end{equation}

de modo que

\begin{equation}
    0\leq\tilde{c}_{(i)}\leq1.
\end{equation}

Como os atributos estão ordenados por preferência,

\begin{equation}
    \tilde{c}_{(1)}
    \geq
    \tilde{c}_{(2)}
    \geq
    \cdots
    \geq
    \tilde{c}_{(n)}.
\end{equation}

A diferença de consenso entre duas posições consecutivas é definida por

\begin{equation}
    \Delta_i
    =
    \tilde{c}_{(i)}
    -
    \tilde{c}_{(i+1)},
    \qquad
    i=1,\ldots,n-1.
    \label{eq:queda-consenso-ac-efs}
\end{equation}

Uma estratégia inicial de corte adaptativo consiste em identificar a maior
separação entre posições consecutivas:

\begin{equation}
    k_{\Delta}
    =
    \underset{1\leq i<n}
    {\arg\max}
    \;
    \Delta_i.
    \label{eq:corte-maior-queda-ac-efs}
\end{equation}

O subconjunto resultante é então

\begin{equation}
    F^{*}_{\Delta}
    =
    \left\{
        f_{(1)},f_{(2)},\ldots,f_{(k_{\Delta})}
    \right\}.
    \label{eq:subset-adaptativo-ac-efs}
\end{equation}

Essa regra interpreta uma queda acentuada no consenso como uma possível
fronteira entre atributos fortemente apoiados pelo ensemble e atributos com
menor suporte relativo. Entretanto, sua efetividade será avaliada
experimentalmente, não sendo assumido \textit{a priori} que a maior diferença
entre escores corresponda necessariamente ao melhor subconjunto para
classificação.


% ----------------------------------------------------------------------------
\subsection{Estratégias de corte consideradas}
\label{subsec:ac-efs-estrategias-corte}
% ----------------------------------------------------------------------------

As estratégias inicialmente consideradas na Etapa 5 podem ser sintetizadas
conforme a Tabela~\ref{tab:cortes-ac-efs}.

\begin{table}[ht]
    \centering
    \small
    \caption{Estratégias de definição do subconjunto final no
    \textit{EFS4AC}.}
    \label{tab:cortes-ac-efs}

    \begin{tabular}{p{2.5cm}p{3.3cm}p{6.5cm}}
        \hline
        \textbf{Estratégia}
        & \textbf{Ponto de corte}
        & \textbf{Finalidade} \\
        \hline

        Top-25\%
        & $k=\lceil0{,}25n\rceil$
        & Avaliar uma redução mais restritiva do conjunto de atributos. \\

        Top-50\%
        & $k=\lceil0{,}50n\rceil$
        & Avaliar uma configuração intermediária de redução. \\

        Top-75\%
        & $k=\lceil0{,}75n\rceil$
        & Avaliar uma redução mais conservadora do espaço de atributos. \\

        Adaptativo
        & $k_{\Delta}=\arg\max_i\Delta_i$
        & Determinar o corte a partir da distribuição dos escores de consenso. \\

        \hline
    \end{tabular}
\end{table}

A comparação dessas configurações permite investigar se um ponto de corte
determinado pela estrutura do consenso é capaz de produzir subconjuntos
competitivos em relação aos percentuais fixos. Ao mesmo tempo, os cortes
percentuais fornecem referências simples e reproduzíveis para avaliar o
comportamento da estratégia adaptativa.


% ----------------------------------------------------------------------------
\subsection{Saída da Etapa 5}
\label{subsec:ac-efs-saida-etapa5}
% ----------------------------------------------------------------------------

Independentemente da estratégia empregada, a saída desta etapa é um ponto de
corte $k^{*}$ e o correspondente subconjunto

\begin{equation}
    F^{*}
    =
    \left\{
        f_{(1)},f_{(2)},\ldots,f_{(k^{*})}
    \right\},
    \qquad
    |F^{*}|=k^{*}.
    \label{eq:subset-final-ac-efs}
\end{equation}

Além do subconjunto selecionado, serão registrados o valor de $k^{*}$, a
proporção de atributos preservados e a proporção de redução obtida:

\begin{equation}
    \operatorname{Redução}(F^{*})
    =
    1-\frac{|F^{*}|}{n}.
    \label{eq:reducao-ac-efs}
\end{equation}

Essas informações permitem analisar separadamente o desempenho preditivo e o
grau de redução promovido pelo método.

O conjunto de dados de treinamento é então projetado sobre $F^{*}$:

\begin{equation}
    D_{\mathrm{tr}}^{*}
    =
    \pi_{F^{*}}
    \left(
        D_{\mathrm{tr}}
    \right),
    \label{eq:projecao-ac-efs}
\end{equation}

e essa representação reduzida é utilizada na etapa subsequente de indução do
classificador associativo.

Assim, a Etapa 5 encerra o processo de seleção realizado pelo
\textit{EFS4AC}: as Etapas 1 e 2 constroem as evidências individuais, a
Etapa 3 analisa sua complementaridade, a Etapa 4 produz o consenso e a
Etapa 5 converte esse consenso em uma decisão efetiva sobre quais atributos
serão preservados.
% ============================================================================
% ============================================================================
\section{Integração do \textit{EFS4AC} com a Classificação Associativa}
\label{sec:ac-efs-integracao}
% ============================================================================

Concluídas as cinco etapas do \textit{EFS4AC}, obtém-se o subconjunto final de
atributos $F^{*}\subseteq F$. A partir desse ponto, o método é integrado ao
processo de Classificação Associativa (CA) como uma etapa supervisionada de
seleção de atributos anterior à indução das regras. Dessa forma, o
\textit{EFS4AC} não modifica diretamente os mecanismos internos do algoritmo
de CA, tais como geração, ordenação ou poda de regras, mas modifica a
representação dos dados sobre a qual esses mecanismos serão executados.

Seja $D_{\mathrm{tr}}$ o conjunto de treinamento originalmente descrito pelo
conjunto de atributos $F$. Após a execução do \textit{EFS4AC}, os dados são
projetados sobre o subconjunto selecionado:

\begin{equation}
    D_{\mathrm{tr}}^{*}
    =
    \pi_{F^{*}}
    \left(
        D_{\mathrm{tr}}
    \right),
    \label{eq:projecao-treino-ca-ac-efs}
\end{equation}

em que $\pi_{F^{*}}(\cdot)$ representa a operação de projeção que preserva
somente os atributos pertencentes a $F^{*}$, além do atributo de classe.

O classificador associativo é então induzido utilizando $D_{\mathrm{tr}}^{*}$:

\begin{equation}
    \mathcal{M}_{CA}
    =
    \mathcal{L}_{CA}
    \left(
        D_{\mathrm{tr}}^{*}
    \right),
    \label{eq:inducao-ca-ac-efs}
\end{equation}

em que $\mathcal{L}_{CA}$ representa o algoritmo de indução do classificador
associativo e $\mathcal{M}_{CA}$ o modelo resultante.

Consequentemente, as regras de associação de classe são extraídas a partir de
uma representação reduzida dos dados. O conjunto $F^{*}$ determina, portanto,
quais atributos poderão participar dos antecedentes das regras induzidas,
enquanto o atributo de classe permanece disponível como consequente durante
o processo de mineração.

A Figura~\ref{fig:posicionamento-ac-efs-ca} apresenta o posicionamento do
\textit{EFS4AC} no fluxo de indução da Classificação Associativa.

\begin{figure}[ht]
    \centering

    \begin{tikzpicture}[
        every node/.style={
            font=\small,
            align=center
        },
        bloco/.style={
            draw,
            rectangle,
            text width=2.25cm,
            minimum height=1.45cm,
            inner sep=4pt
        },
        destaque/.style={
            draw,
            rectangle,
            very thick,
            text width=2.25cm,
            minimum height=1.45cm,
            inner sep=4pt
        },
        seta/.style={
            thick,
            ->,
            >=stealth
        }
    ]

    \node[bloco] (dados) {
        \textbf{Dados de}\\
        treinamento
    };

    \node[bloco, right=0.45cm of dados] (pre) {
        \textbf{Pré-processamento}\\
        e discretização
    };

    \node[destaque, right=0.45cm of pre] (efs) {
        \textbf{\textit{EFS4AC}}\\
        seleção consensual\\
        $F^{*}$
    };

    \node[bloco, right=0.45cm of efs] (regras) {
        \textbf{Indução}\\
        das regras\\
        associativas
    };

    \node[bloco, right=0.45cm of regras] (modelo) {
        \textbf{Classificador}\\
        ordenação, poda\\
        e predição
    };

    \draw[seta] (dados) -- (pre);
    \draw[seta] (pre) -- (efs);
    \draw[seta] (efs) -- (regras);
    \draw[seta] (regras) -- (modelo);

    \end{tikzpicture}

    \caption{Posicionamento do \textit{EFS4AC} no fluxo de indução da
    Classificação Associativa.}
    \label{fig:posicionamento-ac-efs-ca}
\end{figure}

Essa integração permite investigar se a seleção consensual de atributos altera
não apenas o desempenho preditivo, mas também propriedades estruturais do
classificador associativo. Em particular, a redução do conjunto de atributos
pode produzir três efeitos principais: (i) redução do espaço de busca
disponível para a geração de regras; (ii) alteração da quantidade e da
complexidade das regras induzidas; e (iii) alteração do desempenho preditivo
do classificador.

Esses efeitos não são assumidos como consequências necessariamente positivas
do método. A remoção de atributos irrelevantes ou redundantes pode favorecer
a construção do modelo, enquanto a exclusão de atributos informativos pode
prejudicar sua capacidade discriminativa. Por esse motivo, os efeitos do
\textit{EFS4AC} serão avaliados experimentalmente considerando conjuntamente
o desempenho preditivo, o grau de redução de atributos e a complexidade do
modelo associativo resultante.


% ----------------------------------------------------------------------------
\subsection{Cuidados com vazamento de informação}
\label{subsec:ac-efs-vazamento}
% ----------------------------------------------------------------------------

Como os métodos-base utilizados pelo \textit{EFS4AC} empregam informação da
classe para estimar a relevância dos atributos, todo o processo de seleção deve
ser realizado exclusivamente a partir dos dados disponíveis para treinamento.
Nenhuma informação proveniente das instâncias destinadas à avaliação final
pode participar da construção dos rankings, da definição da composição do
ensemble ou da escolha do subconjunto $F^{*}$.

Assim, todas as operações dependentes dos dados devem ser ajustadas utilizando
exclusivamente $D_{\mathrm{tr}}$. Isso inclui o pré-processamento e a
discretização quando seus parâmetros forem estimados a partir dos dados, a
execução dos métodos-base, a geração dos rankings, a análise da diversidade,
a agregação e a determinação do ponto de corte.

Seja $D_{\mathrm{te}}$ o conjunto reservado para avaliação. Após a obtenção de
$F^{*}$ exclusivamente a partir de $D_{\mathrm{tr}}$, o mesmo subconjunto é
utilizado para projetar os dados de teste:

\begin{equation}
    D_{\mathrm{te}}^{*}
    =
    \pi_{F^{*}}
    \left(
        D_{\mathrm{te}}
    \right).
    \label{eq:projecao-teste-ca-ac-efs}
\end{equation}

Dessa maneira, o conjunto de teste não participa da decisão sobre quais
atributos serão selecionados. Sua função restringe-se à avaliação do modelo
induzido a partir dos dados de treinamento.


% ----------------------------------------------------------------------------
\subsection{Execução em procedimentos de validação cruzada}
\label{subsec:ac-efs-validacao-cruzada}
% ----------------------------------------------------------------------------

Em procedimentos de validação cruzada, essa separação deve ser preservada em
cada partição. Para um procedimento com $K$ \textit{folds}, seja
$D_{\mathrm{tr}}^{(q)}$ o conjunto de treinamento e
$D_{\mathrm{va}}^{(q)}$ o conjunto de validação correspondente ao
\textit{fold} $q$.

Para cada $q=1,\ldots,K$, o \textit{EFS4AC} deve ser executado novamente sobre
$D_{\mathrm{tr}}^{(q)}$:

\begin{equation}
    F_{q}^{*}
    =
    \operatorname{EFS4AC}
    \left(
        D_{\mathrm{tr}}^{(q)}
    \right).
    \label{eq:efs-fold-ac-efs}
\end{equation}

Em seguida, treinamento e validação são projetados utilizando exclusivamente
o subconjunto obtido naquele \textit{fold}:

\begin{equation}
    D_{\mathrm{tr}}^{*(q)}
    =
    \pi_{F_q^{*}}
    \left(
        D_{\mathrm{tr}}^{(q)}
    \right),
\end{equation}

\begin{equation}
    D_{\mathrm{va}}^{*(q)}
    =
    \pi_{F_q^{*}}
    \left(
        D_{\mathrm{va}}^{(q)}
    \right).
\end{equation}

O classificador associativo é então induzido em
$D_{\mathrm{tr}}^{*(q)}$ e avaliado em $D_{\mathrm{va}}^{*(q)}$.

Esse procedimento implica que diferentes \textit{folds} podem produzir
diferentes rankings consensuais e diferentes subconjuntos $F_q^{*}$. Tal
variação não constitui uma inconsistência do protocolo, mas uma consequência
natural de estimar a seleção de atributos exclusivamente a partir dos dados
disponíveis em cada conjunto de treinamento.

A execução integral do \textit{EFS4AC} dentro de cada partição garante que a
estimativa de desempenho do classificador associativo reflita todo o processo
de seleção e indução, evitando que informações das instâncias de validação ou
teste influenciem antecipadamente a escolha dos atributos.
% ============================================================================
% ============================================================================
\section{Procedimento Geral do Método}
\label{sec:ac-efs-procedimento}
% ============================================================================

As etapas anteriormente descritas constituem um fluxo sequencial no qual a
saída de cada etapa fornece as informações necessárias para a etapa seguinte.
De forma geral, o \textit{EFS4AC} recebe um conjunto de treinamento
$D_{\mathrm{tr}}$ e uma configuração experimental definida pelos métodos-base,
pela estratégia de controle da diversidade, pelo operador de agregação e pela
regra de corte. Como resultado, produz um ranking consensual dos atributos e
um subconjunto final $F^{*}$.

O procedimento completo pode ser sintetizado da seguinte forma:

\begin{enumerate}

    \item definir o conjunto de métodos-base
    $\mathcal{S}=\{S_1,S_2,\ldots,S_M\}$ que compõe o ensemble, de acordo com
    os critérios estabelecidos na Etapa 1;

    \item aplicar cada seletor $S_j\in\mathcal{S}$ exclusivamente ao conjunto
    de treinamento $D_{\mathrm{tr}}$;

    \item converter as saídas dos seletores em rankings completos
    $R_j$ e organizá-los no conjunto
    $\mathcal{R}=\{R_1,R_2,\ldots,R_M\}$ e na matriz de rankings
    $\mathbf{R}$;

    \item calcular a similaridade entre os rankings, construir a matriz
    $\mathbf{S}$ e aplicar a estratégia de controle da redundância definida
    para a Etapa 3;

    \item obter, conforme a configuração adotada, o conjunto de rankings
    disponibilizado para agregação $\mathcal{R}^{*}$ e, quando aplicável,
    informações adicionais de agrupamento ou ponderação;

    \item aplicar o operador de agregação $\mathcal{A}$ aos rankings de
    $\mathcal{R}^{*}$, produzindo os escores de consenso
    $\mathbf{c}$ e o ranking consensual $R^{*}$;

    \item aplicar a função de corte $\mathcal{K}$ sobre o ranking consensual e
    seus respectivos escores, determinando o ponto de corte $k^{*}$ e o
    subconjunto final de atributos $F^{*}$;

    \item projetar o conjunto de treinamento sobre os atributos selecionados,
    obtendo $D_{\mathrm{tr}}^{*}=\pi_{F^{*}}(D_{\mathrm{tr}})$;

    \item induzir o classificador associativo utilizando
    $D_{\mathrm{tr}}^{*}$;

    \item registrar as saídas finais e os diagnósticos intermediários
    necessários à análise experimental, incluindo os rankings individuais,
    a matriz de similaridade, o ranking consensual, os escores de consenso,
    o ponto de corte e o subconjunto final selecionado.

\end{enumerate}

As cinco primeiras transformações metodológicas do \textit{EFS4AC} podem ser
expressas, de forma compacta, pela composição

\begin{equation}
    D_{\mathrm{tr}}
    \xrightarrow{\mathcal{S}}
    \mathbf{R}
    \xrightarrow{\mathcal{D}}
    \mathcal{R}^{*}
    \xrightarrow{\mathcal{A}}
    \left(R^{*},\mathbf{c}\right)
    \xrightarrow{\mathcal{K}}
    F^{*},
    \label{eq:fluxo-geral-ac-efs}
\end{equation}

em que $\mathcal{S}$ representa a aplicação dos métodos-base,
$\mathcal{D}$ o mecanismo de análise e controle da diversidade,
$\mathcal{A}$ o operador de agregação e $\mathcal{K}$ a regra responsável
pela definição do subconjunto final.

A Tabela~\ref{tab:entradas-saidas-ac-efs} sintetiza as entradas, operações e
saídas correspondentes às cinco etapas centrais do método.

\begin{table}[ht]
    \centering
    \small
    \caption{Síntese das etapas do método \textit{EFS4AC}.}
    \label{tab:entradas-saidas-ac-efs}
    \renewcommand{\arraystretch}{1.20}
    \setlength{\tabcolsep}{4pt}

    \begin{tabular}{
        p{0.7cm}
        p{3.0cm}
        p{3.0cm}
        p{3.3cm}
        p{3.2cm}
    }
        \hline
        \textbf{Et.}
        & \textbf{Objetivo}
        & \textbf{Entrada}
        & \textbf{Operação principal}
        & \textbf{Saída} \\
        \hline

        1
        & Definir os métodos-base
        & Métodos candidatos e critérios de inclusão
        & Composição do ensemble
        & Conjunto de seletores $\mathcal{S}$ \\

        2
        & Padronizar as saídas
        & $D_{\mathrm{tr}}$ e $\mathcal{S}$
        & Geração de rankings completos
        & $\mathcal{R}$ e matriz $\mathbf{R}$ \\

        3
        & Analisar e controlar redundância
        & Matriz $\mathbf{R}$
        & Similaridade e controle de diversidade
        & $\mathbf{S}$, $\mathcal{R}^{*}$ e estrutura de controle \\

        4
        & Construir o consenso
        & $\mathcal{R}^{*}$ e configuração de agregação
        & Borda, RRF ou TOPSIS
        & $R^{*}$ e escores $\mathbf{c}$ \\

        5
        & Definir o subconjunto final
        & $R^{*}$ e $\mathbf{c}$
        & Aplicação da regra de corte
        & $k^{*}$ e subconjunto $F^{*}$ \\

        \hline
    \end{tabular}
\end{table}

Ao término dessas cinco etapas, o \textit{EFS4AC} fornece uma decisão
estruturada sobre a relevância e a seleção dos atributos. O ranking
$R^{*}$ preserva a ordenação consensual produzida pelo ensemble, enquanto
$F^{*}$ materializa essa decisão na forma de um subconjunto de atributos.

A projeção dos dados e a posterior indução do classificador associativo não
constituem novas etapas internas do \textit{EFS4AC}, mas correspondem à sua
integração com o processo de Classificação Associativa, conforme apresentado
na Seção~\ref{sec:ac-efs-integracao}.

% ============================================================================
% ============================================================================
\section{Configurações Metodológicas a Serem Investigadas}
\label{sec:ac-efs-configuracoes}
% ============================================================================

A arquitetura do \textit{EFS4AC} foi definida de forma modular para permitir
que decisões correspondentes a diferentes etapas do método sejam analisadas
separadamente. Dessa forma, além de avaliar o desempenho final do método,
pretende-se investigar como a composição do ensemble, o controle da
redundância, o mecanismo de agregação e a regra de corte influenciam os
subconjuntos produzidos e, posteriormente, o comportamento do classificador
associativo.

As configurações experimentais serão organizadas em quatro fatores principais:
(i) composição do ensemble; (ii) estratégia de controle da diversidade;
(iii) operador de agregação; e (iv) regra de definição do subconjunto final.
Além desses fatores, serão consideradas configurações de referência externas
ao ensemble para fins de comparação.


% ----------------------------------------------------------------------------
\subsection{Composição do ensemble}
\label{subsec:ac-efs-config-composicao}
% ----------------------------------------------------------------------------

O primeiro fator investiga a influência da natureza dos métodos-base sobre o
consenso produzido pelo \textit{EFS4AC}. Serão consideradas inicialmente três
composições:

\begin{equation}
    E_U
    =
    \{
        \mathrm{Gini},
        \chi^2,
        \mathrm{SU},
        \mathrm{MIM}
    \},
\end{equation}

\begin{equation}
    E_M
    =
    \{
        \mathrm{MRMR},
        \mathrm{JMI},
        \mathrm{CIFE},
        \mathrm{CMIM},
        \mathrm{ICAP},
        \mathrm{MRI}
    \},
\end{equation}

e

\begin{equation}
    E_{U+M}
    =
    E_U\cup E_M.
\end{equation}

A configuração $E_U$ reúne métodos univariados, enquanto $E_M$ reúne métodos
multivariados/sequenciais. A configuração $E_{U+M}$ combina as duas famílias
e representa o cenário de maior heterogeneidade metodológica entre os
seletores inicialmente considerados.

A comparação dessas configurações permitirá analisar se ensembles constituídos
por métodos de uma mesma família e ensembles heterogêneos apresentam
comportamentos distintos quanto aos rankings consensuais e aos subconjuntos
selecionados.

Entretanto, essa comparação deve ser interpretada considerando que as
configurações possuem números diferentes de componentes: $E_U$, $E_M$ e
$E_{U+M}$ contêm, respectivamente, 4, 6 e 10 seletores. Assim, eventuais
diferenças observadas não podem ser atribuídas exclusivamente à
heterogeneidade, uma vez que a quantidade de métodos-base também varia entre
as configurações.


% ----------------------------------------------------------------------------
\subsection{Estratégia de controle da diversidade}
\label{subsec:ac-efs-config-diversidade}
% ----------------------------------------------------------------------------

O segundo fator corresponde ao tratamento da similaridade entre os rankings
produzidos pelos métodos-base. Conforme definido na Etapa 3, serão investigadas
as seguintes configurações:

\begin{itemize}

    \item \textbf{$D_0$ -- sem controle:}
    todos os rankings produzidos pelos métodos-base são disponibilizados
    diretamente para a agregação;

    \item \textbf{$D_L$ -- controle por limiar:}
    a matriz de similaridade é utilizada para identificar rankings com
    concordância elevada e controlar sua influência no ensemble;

    \item \textbf{$D_C$ -- controle por agrupamento:}
    os rankers são organizados segundo sua similaridade, permitindo controlar
    a influência de grupos formados por rankings altamente concordantes.

\end{itemize}

A configuração $D_0$ funciona como referência interna para as estratégias
$D_L$ e $D_C$, permitindo verificar se o controle explícito da redundância
entre rankings modifica de maneira relevante o consenso produzido pelo
ensemble.

Quando uma estratégia da Etapa 3 resultar em pesos associados aos rankers ou
aos grupos identificados, esses valores serão encaminhados à Etapa 4 e
incorporados pelos respectivos operadores de agregação.


% ----------------------------------------------------------------------------
\subsection{Operador de agregação}
\label{subsec:ac-efs-config-agregacao}
% ----------------------------------------------------------------------------

O terceiro fator corresponde ao mecanismo utilizado para transformar os
rankings individuais em um ranking consensual. Serão comparados três
operadores:

\begin{itemize}

    \item \textbf{Borda}, baseado na soma das pontuações derivadas das posições
    dos atributos;

    \item \textbf{\textit{Reciprocal Rank Fusion} (RRF)}, que atribui maior
    influência relativa às primeiras posições dos rankings;

    \item \textbf{TOPSIS}, que formula a agregação como um problema de decisão
    multicritério e ordena os atributos de acordo com sua proximidade à
    solução ideal.

\end{itemize}

Os três operadores receberão, para cada configuração experimental, as mesmas
informações provenientes da Etapa 3. Dessa forma, busca-se analisar o efeito
do mecanismo de agregação sem modificar simultaneamente a composição dos
rankings fornecidos como entrada.


% ----------------------------------------------------------------------------
\subsection{Regra de definição do subconjunto}
\label{subsec:ac-efs-config-corte}
% ----------------------------------------------------------------------------

O quarto fator corresponde à transformação do ranking consensual em um
subconjunto efetivo de atributos. Serão consideradas quatro configurações:

\begin{equation}
    \mathcal{K}
    \in
    \{
        K_{25},
        K_{50},
        K_{75},
        K_{\Delta}
    \},
\end{equation}

em que $K_{25}$, $K_{50}$ e $K_{75}$ correspondem, respectivamente, aos
cortes Top-25\%, Top-50\% e Top-75\%, enquanto $K_{\Delta}$ representa a
estratégia adaptativa baseada na maior queda entre escores consecutivos de
consenso, conforme definida na Etapa 5.

Os cortes percentuais fornecem níveis padronizados de redução, enquanto
$K_{\Delta}$ permite investigar se a própria estrutura do consenso pode ser
utilizada para determinar automaticamente a quantidade de atributos
selecionados.


% ----------------------------------------------------------------------------
\subsection{Configurações de referência}
\label{subsec:ac-efs-config-referencia}
% ----------------------------------------------------------------------------

Além das configurações internas do \textit{EFS4AC}, serão considerados
cenários de referência para contextualizar seus resultados. O primeiro
corresponde à indução do classificador associativo utilizando o conjunto
completo de atributos, isto é, sem aplicação de seleção de atributos.

O segundo grupo de referências corresponde à utilização individual dos
métodos-base de FS. Nesse cenário, os seletores são avaliados separadamente,
sem agregação de rankings, permitindo comparar os subconjuntos obtidos pelo
ensemble com aqueles produzidos pelos seus componentes individuais.

Essas referências permitem distinguir os efeitos associados à redução
dimensional daqueles decorrentes especificamente da construção de um consenso
entre diferentes métodos de seleção.


% ----------------------------------------------------------------------------
\subsection{Síntese das configurações}
\label{subsec:ac-efs-config-sintese}
% ----------------------------------------------------------------------------

A Tabela~\ref{tab:configuracoes-ac-efs} sintetiza os fatores metodológicos
previstos para a investigação experimental.

\begin{table}[ht]
    \centering
    \small
    \caption{Fatores e configurações metodológicas previstas para o
    \textit{EFS4AC}.}
    \label{tab:configuracoes-ac-efs}
    \renewcommand{\arraystretch}{1.20}

    \begin{tabular}{p{3.3cm}p{4.8cm}p{5.3cm}}
        \hline
        \textbf{Fator}
        & \textbf{Configurações}
        & \textbf{Questão investigada} \\
        \hline

        Composição do ensemble
        & $E_U$, $E_M$, $E_{U+M}$
        & Qual o efeito da composição dos métodos-base sobre o consenso? \\

        Controle da diversidade
        & $D_0$, $D_L$, $D_C$
        & O controle de rankings redundantes modifica o resultado do ensemble? \\

        Agregação
        & Borda, RRF, TOPSIS
        & Como diferentes mecanismos de consenso afetam a ordenação dos atributos? \\

        Regra de corte
        & $K_{25}$, $K_{50}$, $K_{75}$, $K_{\Delta}$
        & Como a regra de corte influencia o subconjunto final selecionado? \\

        Referências
        & Sem FS e métodos individuais
        & O ensemble apresenta comportamento distinto de suas referências? \\

        \hline
    \end{tabular}
\end{table}

A análise dessas configurações permitirá investigar separadamente as principais
decisões que compõem o \textit{EFS4AC}. Em particular, pretende-se verificar
como a composição dos métodos-base influencia o consenso, se o controle de
redundância entre rankings altera os resultados da agregação, como diferentes
operadores produzem ordenações consensuais distintas e se uma estratégia
adaptativa de corte apresenta comportamento competitivo em relação aos
percentuais fixos.

Essas comparações serão realizadas sob um mesmo protocolo experimental,
mantendo constantes, sempre que possível, os demais componentes do método ao
analisar cada fator. Dessa forma, busca-se reduzir a atribuição de diferenças
observadas a alterações simultâneas em múltiplos componentes da arquitetura.
% ============================================================================
% ============================================================================
\section{Saídas e Critérios de Análise}
\label{sec:ac-efs-saidas}
% ============================================================================

A avaliação do \textit{EFS4AC} não será restrita ao desempenho preditivo do
classificador associativo. Como o método envolve diferentes decisões
intermediárias --- composição do ensemble, controle da diversidade, agregação
e definição do ponto de corte ---, serão registrados tanto os resultados
finais quanto informações produzidas ao longo de suas etapas.

Esse registro permitirá analisar não apenas se determinado subconjunto produz
bom desempenho de classificação, mas também como ele foi obtido, qual redução
dimensional foi alcançada, quais efeitos foram produzidos sobre a estrutura do
classificador associativo e quão consistente é a seleção diante de variações
nos dados de treinamento.

Entre as principais saídas intermediárias e finais a serem registradas,
destacam-se:

\begin{itemize}

    \item os rankings individuais $R_j$ produzidos pelos métodos-base;

    \item a matriz de rankings $\mathbf{R}$;

    \item a matriz de similaridade $\mathbf{S}$ entre os rankers;

    \item a composição do ensemble após a aplicação da estratégia de controle
    da diversidade e, quando aplicável, os grupos ou pesos associados aos
    rankers;

    \item os escores de consenso $\mathbf{c}$ e o ranking consensual $R^{*}$;

    \item o ponto de corte $k^{*}$ e a regra de corte utilizada;

    \item o subconjunto final de atributos $F^{*}$;

    \item o tempo de execução das etapas do \textit{EFS4AC};

    \item o tempo de indução do classificador associativo a partir do conjunto
    de atributos selecionado.

\end{itemize}

A análise das diferentes configurações do método será organizada em quatro
dimensões principais: desempenho preditivo, redução dimensional, complexidade
do classificador associativo e estabilidade da seleção.


% ----------------------------------------------------------------------------
\subsection{Desempenho preditivo}
\label{subsec:ac-efs-avaliacao-preditiva}
% ----------------------------------------------------------------------------

A primeira dimensão corresponde à capacidade preditiva do classificador
associativo induzido após a seleção de atributos. A medida principal será a
F1-Macro, calculada a partir das predições realizadas sobre dados não
utilizados na construção do subconjunto $F^{*}$.

A utilização da F1-Macro permite atribuir a mesma importância às diferentes
classes no cálculo da medida global, independentemente da frequência relativa
de cada uma delas. Para um problema com $C$ classes, define-se

\begin{equation}
    F1_{\mathrm{Macro}}
    =
    \frac{1}{C}
    \sum_{c=1}^{C}
    F1_c,
    \label{eq:f1-macro-ac-efs}
\end{equation}

em que

\begin{equation}
    F1_c
    =
    2
    \frac{
        \operatorname{Precisão}_c
        \operatorname{Revocação}_c
    }{
        \operatorname{Precisão}_c+
        \operatorname{Revocação}_c
    }.
\end{equation}

Os resultados obtidos pelas diferentes configurações do \textit{EFS4AC} serão
comparados com os cenários de referência definidos anteriormente, incluindo o
classificador associativo induzido sem seleção de atributos e os resultados
obtidos pelos métodos individuais de FS.

Essa comparação permitirá analisar se as alterações observadas no desempenho
estão associadas apenas à redução dimensional ou se o consenso entre múltiplos
seletores apresenta comportamento distinto das seleções individuais.


% ----------------------------------------------------------------------------
\subsection{Redução dimensional}
\label{subsec:ac-efs-avaliacao-reducao}
% ----------------------------------------------------------------------------

A segunda dimensão corresponde ao grau de redução do espaço original de
atributos. Para cada subconjunto $F^{*}$ serão registrados tanto o número
absoluto de atributos selecionados, $|F^{*}|$, quanto sua proporção em relação
ao conjunto original $F$.

A taxa de redução é definida por

\begin{equation}
    \operatorname{Red}(F^{*})
    =
    1-
    \frac{|F^{*}|}{|F|}.
    \label{eq:reducao-atributos-ac-efs}
\end{equation}

Valores próximos de zero indicam pouca redução do conjunto original, enquanto
valores maiores representam subconjuntos proporcionalmente menores.

A taxa de redução não será interpretada isoladamente como indicadora de
qualidade. Uma redução elevada pode ser desejável do ponto de vista de
dimensionalidade, mas pode também eliminar atributos necessários à
classificação. Por essa razão, $\operatorname{Red}(F^{*})$ será analisada em
conjunto com o desempenho preditivo e com a complexidade do modelo associativo
resultante.


% ----------------------------------------------------------------------------
\subsection{Complexidade do classificador associativo}
\label{subsec:ac-efs-avaliacao-complexidade}
% ----------------------------------------------------------------------------

A terceira dimensão busca analisar como a seleção de atributos realizada pelo
\textit{EFS4AC} afeta a estrutura do classificador associativo induzido.

Como a redução do conjunto de atributos modifica o espaço disponível para a
formação dos antecedentes das regras, serão registrados, inicialmente:

\begin{itemize}

    \item o número total de regras do classificador final;

    \item o tamanho médio dos antecedentes das regras;

    \item o tempo necessário para a indução do classificador associativo.

\end{itemize}

Seja $\mathcal{Q}$ o conjunto final de regras do classificador associativo e
seja $\operatorname{ant}(q)$ o antecedente de uma regra $q\in\mathcal{Q}$.
O tamanho médio dos antecedentes pode ser expresso por

\begin{equation}
    \overline{L}_{\mathrm{ant}}
    =
    \frac{1}{|\mathcal{Q}|}
    \sum_{q\in\mathcal{Q}}
    |\operatorname{ant}(q)|.
    \label{eq:tamanho-medio-regras-ac-efs}
\end{equation}

Essas medidas permitirão verificar se diferentes configurações de seleção
produzem alterações na quantidade e no tamanho das regras, além de mudanças no
custo computacional de indução.

Assim como ocorre com a redução dimensional, menor complexidade não será
interpretada isoladamente como superioridade. A análise considerará
conjuntamente a estrutura do modelo e seu desempenho preditivo.


% ----------------------------------------------------------------------------
\subsection{Estabilidade da seleção}
\label{subsec:ac-efs-avaliacao-estabilidade}
% ----------------------------------------------------------------------------

A quarta dimensão corresponde à estabilidade dos subconjuntos produzidos pelo
\textit{EFS4AC}. Neste trabalho, estabilidade refere-se à consistência da
seleção quando o método é aplicado a diferentes partições dos dados de
treinamento.

Sejam

\begin{equation}
    F^{*}_{1},F^{*}_{2},\ldots,F^{*}_{K}
\end{equation}

os subconjuntos obtidos em $K$ partições de treinamento. A similaridade entre
dois subconjuntos $F^{*}_{a}$ e $F^{*}_{b}$ pode ser calculada pelo índice de
Jaccard:

\begin{equation}
    J
    \left(
        F^{*}_{a},F^{*}_{b}
    \right)
    =
    \frac{
        |F^{*}_{a}\cap F^{*}_{b}|
    }{
        |F^{*}_{a}\cup F^{*}_{b}|
    }.
    \label{eq:jaccard-estabilidade-ac-efs}
\end{equation}

A estabilidade global pode então ser expressa pela média das similaridades
entre todos os pares de subconjuntos:

\begin{equation}
    \operatorname{Estab}
    =
    \frac{2}{K(K-1)}
    \sum_{a=1}^{K-1}
    \sum_{b=a+1}^{K}
    J
    \left(
        F^{*}_{a},F^{*}_{b}
    \right).
    \label{eq:estabilidade-media-ac-efs}
\end{equation}

Valores mais próximos de 1 indicam maior sobreposição entre os subconjuntos
produzidos nas diferentes partições, enquanto valores menores indicam maior
variabilidade na composição de $F^{*}$.

A estabilidade será analisada em conjunto com a quantidade de atributos
selecionados, uma vez que a sobreposição entre subconjuntos também pode ser
influenciada por seus tamanhos. Dessa forma, essa medida será utilizada como
indicador complementar da consistência do processo de seleção, e não como
critério isolado para determinar a melhor configuração.


% ----------------------------------------------------------------------------
\subsection{Custo computacional}
\label{subsec:ac-efs-avaliacao-tempo}
% ----------------------------------------------------------------------------

Embora o tempo de execução esteja relacionado à complexidade computacional,
serão registrados separadamente o custo da seleção e o custo da indução do
classificador. Para cada execução, serão considerados:

\begin{equation}
    T_{\mathrm{EFS}}
\end{equation}

como o tempo total necessário para executar as etapas do \textit{EFS4AC}, e

\begin{equation}
    T_{\mathrm{CA}}
\end{equation}

como o tempo necessário para induzir o classificador associativo após a
seleção.

Quando pertinente, também poderá ser analisado o tempo total do fluxo:

\begin{equation}
    T_{\mathrm{total}}
    =
    T_{\mathrm{EFS}}
    +
    T_{\mathrm{CA}}.
    \label{eq:tempo-total-ac-efs}
\end{equation}

Essa separação é importante porque uma configuração pode introduzir maior custo
na etapa de seleção e, simultaneamente, reduzir o custo posterior de indução do
classificador associativo.


% ----------------------------------------------------------------------------
\subsection{Análise conjunta dos critérios}
\label{subsec:ac-efs-avaliacao-conjunta}
% ----------------------------------------------------------------------------

Os critérios anteriores serão analisados de forma complementar. O objetivo não
é assumir que uma única medida seja suficiente para caracterizar o
\textit{EFS4AC}, mas observar os efeitos das diferentes configurações sobre
predição, dimensionalidade, estrutura do classificador e estabilidade.

A comparação experimental considerará, portanto, o conjunto de resultados

\begin{equation}
    \mathcal{E}
    =
    \left\{
        F1_{\mathrm{Macro}},
        \operatorname{Red},
        |\mathcal{Q}|,
        \overline{L}_{\mathrm{ant}},
        \operatorname{Estab},
        T_{\mathrm{EFS}},
        T_{\mathrm{CA}}
    \right\}.
    \label{eq:criterios-avaliacao-ac-efs}
\end{equation}

Essa análise permitirá caracterizar possíveis relações de compromisso entre
os critérios. Por exemplo, configurações com diferentes graus de redução podem
apresentar comportamentos distintos quanto ao desempenho preditivo, à
estabilidade e à complexidade do classificador associativo. Tais relações
serão examinadas experimentalmente, sem assumir previamente que a otimização
de um único critério implique melhora nos demais.

% ============================================================================
% ============================================================================
\section{Considerações Finais do Capítulo}
\label{sec:ac-efs-consideracoes}
% ============================================================================

Neste capítulo foi apresentado o \textit{EFS4AC}, método de
\textit{Ensemble Feature Selection} proposto para atuar previamente à indução
de Classificadores Associativos. A proposta parte do princípio de que
diferentes métodos de seleção de atributos podem produzir evidências distintas
sobre a relevância dos atributos e, portanto, busca integrar essas evidências
antes da etapa de mineração das regras.

O método foi estruturado em cinco etapas interdependentes. A primeira define
a composição do ensemble a partir de diferentes métodos-base de seleção de
atributos. A segunda converte as saídas desses métodos para uma representação
ordinal comum, produzindo rankings comparáveis. A terceira analisa a
similaridade entre esses rankings e permite controlar a influência de
evidências redundantes. A quarta integra os rankings por meio de operadores de
agregação, produzindo uma ordenação consensual dos atributos. Por fim, a quinta
etapa transforma esse consenso em um subconjunto efetivo de atributos por meio
de uma regra de corte.

De forma sintética, o fluxo central do método pode ser representado por

\begin{equation}
    D_{\mathrm{tr}}
    \xrightarrow{\mathcal{S}}
    \mathbf{R}
    \xrightarrow{\mathcal{D}}
    \mathcal{R}^{*}
    \xrightarrow{\mathcal{A}}
    \left(R^{*},\mathbf{c}\right)
    \xrightarrow{\mathcal{K}}
    F^{*},
    \label{eq:sintese-final-ac-efs}
\end{equation}

em que $\mathcal{S}$ representa a aplicação dos métodos-base,
$\mathcal{D}$ o mecanismo de análise e controle da diversidade,
$\mathcal{A}$ o operador de agregação e $\mathcal{K}$ a regra de corte. O
resultado desse processo é o subconjunto $F^{*}$, utilizado para projetar os
dados antes da indução do classificador associativo.

Uma característica importante da proposta é a separação explícita entre as
decisões envolvidas no processo de seleção. A composição do ensemble, o
controle da redundância entre rankers, a construção do consenso e a definição
do número de atributos selecionados são tratados como problemas distintos.
Essa organização permite analisar individualmente a contribuição de cada
componente e evita associar o resultado final a uma única decisão metodológica.

Além disso, o \textit{EFS4AC} foi concebido de forma modular. Diferentes
composições de métodos-base, estratégias de controle da diversidade,
operadores de agregação e regras de corte podem ser instanciados sem modificar
a estrutura geral do método. Neste trabalho, essa característica será
explorada por meio das composições $E_U$, $E_M$ e $E_{U+M}$, das estratégias
de diversidade definidas anteriormente, dos agregadores Borda, RRF e TOPSIS e
das regras de corte percentuais e adaptativa.

A integração com a Classificação Associativa ocorre externamente ao mecanismo
de indução das regras. O \textit{EFS4AC} determina quais atributos serão
disponibilizados ao classificador, mas não modifica diretamente os
procedimentos internos de geração, ordenação ou poda das regras. Essa
separação permite analisar os efeitos da seleção consensual sobre o processo
associativo mantendo inalterado o algoritmo de classificação utilizado como
referência.

Do ponto de vista experimental, a proposta também não será caracterizada
exclusivamente pelo desempenho preditivo. Sua avaliação considerará
conjuntamente a capacidade preditiva do classificador associativo, a redução
do espaço de atributos, a complexidade estrutural do modelo induzido, a
estabilidade dos subconjuntos selecionados e o custo computacional. Dessa
forma, será possível investigar não apenas se determinada configuração produz
bons resultados de classificação, mas também quais efeitos a seleção exerce
sobre a dimensionalidade dos dados e sobre a estrutura do classificador
baseado em regras.

Outro aspecto fundamental do método é que todas as decisões dependentes dos
dados são realizadas exclusivamente sobre as partições de treinamento. Em
procedimentos de validação cruzada, o \textit{EFS4AC} deve ser executado
integralmente em cada \textit{fold}, incluindo a geração dos rankings, a
análise da diversidade, a agregação e a determinação do subconjunto final.
As partições destinadas à validação ou ao teste são apenas projetadas sobre
os atributos previamente selecionados. Esse procedimento preserva a
independência da avaliação e evita que informações externas ao treinamento
influenciem a seleção dos atributos.

A contribuição metodológica pretendida com o \textit{EFS4AC} consiste,
portanto, em estabelecer uma estrutura sistemática para investigar a
combinação de múltiplas evidências de seleção de atributos no contexto da
Classificação Associativa. Em vez de pressupor que uma determinada composição,
estratégia de diversidade, agregador ou regra de corte seja superior, a
arquitetura proposta permite que essas decisões sejam avaliadas
experimentalmente sob um protocolo comum.

No capítulo seguinte, essa arquitetura será instanciada experimentalmente. As
diferentes configurações do \textit{EFS4AC} serão comparadas entre si e com
cenários de referência, incluindo a utilização do conjunto completo de
atributos e a aplicação individual dos métodos-base. A análise permitirá
examinar as relações entre consenso, redução dimensional, estabilidade,
desempenho preditivo, complexidade das regras e custo computacional,
fornecendo evidências sobre o comportamento da abordagem proposta no contexto
dos Classificadores Associativos.
